\begin{problem}[2020第37届全国中学生物理竞赛复赛][光纤陀螺仪Sagnac效应分析]{122}
    光纤陀螺仪是一种能够精确测定运动物体方位的光学仪器，它是现代航海、航空和航天等领域中被广泛使用的一种惯性导航仪器。光纤陀螺仪导航主要基于下述效应：在一个半径为 $R$、以角速度 $\Omega$ 转动的光纤环路中，从固定在环上的分束器 $A$ 分出的两束相干光分别沿逆时针（CCW）和顺时针（CW）方向传播后回到 $A$，两者的光程不一样。检测两束光的相位差或干涉条纹的变化，可确定该光纤环路的转动角速度 $\Omega$。真空中的光速为 $c$。
    \begin{subproblem}
        光纤由内、外两层介质构成，内、外层介质的折射率分别为 $n_{1}$、$n_{2}$ ($n_{1} > n_{2}$)。为了使光能在光纤内传输，光在输入端口（端口外介质的折射率为 $n_{0}$）的入射角 $i$ 应满足什么条件？
    \end{subproblem}
    \begin{subproblem}
        考虑沿环路密绕 $N$ 匝光纤，首尾相接于分束器 $A$，光纤环路以角速度 $\Omega$ 沿逆时针方向旋转。从 $A$ 分出两束相干光沿光纤环路逆时针和顺时针方向传播，又回到 $A$。已知光传播介质的折射率为 $n_{1}$，两束光在真空中的波长为 $\lambda_{0}$，问两束相干光分别沿光纤环路逆时针和顺时针方向传播的时间 $t_{\mathrm{CCW}}$ 和 $t_{\mathrm{CW}}$ 为多少？这两束光的相位差为多少？试指出该相差与介质折射率 $n_{1}$ 之间的关系。
    \end{subproblem}
    \begin{subproblem}
        为了提高陀螺系统的微型化程度，人们提出了谐振式光学陀螺系统。该系统中含有一个沿逆时针方向旋转的光学环形腔，其半径为 $R$，旋转角速度为 $\Omega$ ($R\Omega << c$)，已知腔中的介质折射率为 $n$，在腔内存在沿逆时针和顺时针方向传播的两类共振模式。当环形腔静止时，这些模式的波长 $\lambda_{\mathrm{m0}}$ 由周期边界条件 $2\pi R = m \lambda_{\mathrm{m0}}$ 决定，其中 $m$ 为正整数，称为共振模式的级次。当环形腔旋转时，同一级次（均为 $m$）的沿顺时针和逆时针方向传播的共振模式存在一个共振频率差 $\Delta f$，试导出环形腔转动角速度 $\Omega$ 与该共振频率差 $\Delta f$ 之间的关系。
    \end{subproblem}
\end{problem}